\bookchapter{Bab 15: Sains Bateri \& Pengoptimuman Ketahanan}

\chapterintro{Bateri ialah “jantung tenaga” telefon pintar anda. Bab ini membantu anda memahami perangai bateri, belajar teknik cas dan ketahanan secara saintifik, dan mengucapkan selamat tinggal kepada kebimbangan paras bateri.}

\booksection{Apakah itu Bateri? Kitaran dan Tahap Kesihatan}

: Suhu tinggi adalah pembunuh utama bateri.

\begin{itemize}
  \item \textbf{Akan menua}: Selepas lama digunakan, kapasiti maksimum akan menurun, ini adalah fenomena biasa.
  \item \textbf{Penunjuk Teras: Bilangan Kitaran dan Tahap Kesihatan}Satu kitaran bukan bermaksud mengecas sekali. Contohnya, hari ini guna 60\% dan cas penuh, esok guna 40\% dan cas penuh, jumlah penggunaan 100\% barulah dikira satu kitaran.
  \item \textbf{Tahap kesihatan bateri adalah peratusan kapasiti maksimum semasa berbanding bateri baru. Apabila di bawah 80\%, anda akan jelas berasa ketahanan bateri semakin pendek.}Bagaimana melihatnya?
\end{itemize}

\booksubsection{Penunjuk Teras: Bilangan Kitaran dan Tahap Kesihatan}

\begin{itemize}
  \item 2. Mitos Pengecasan? Cas Pantas, Cas Wayarles \& Peraturan Emas
  \item Cas pantas merosakkan bateri?
  \item \textbf{Menggunakan aksesori cas pantas ori atau bersijil, kerosakannya sangat kecil. Telefon pintar dan pengecas mempunyai kawalan suhu, dan akan mengurangkan kuasa secara automatik apabila terlalu panas.}Pengecasan wayarles
\end{itemize}

\booksection{Mitos Pengecasan? Cas Pantas, Cas Wayarles \& Peraturan Emas}

\begin{itemize}
  \item \textbf{Peraturan emas penjagaan bateri}: Cuba kekalkan paras bateri antara 20\%-80\%, jangan biarkan ia kelaparan (0\%) atau terlalu kenyang (100\%).
  \item \textbf{Pengecasan wayarles}Anda mungkin pernah mengalami situasi ini: telefon pintar anda menyokong cas pantas, tetapi selepas menukar pengecas, ia menjadi sangat perlahan. Ini kerana cas pantas bukan sekadar "kuasa besar sudah cukup", tetapi memerlukan pengecas, kabel data, dan telefon pintar bertiga "bercakap dalam bahasa yang sama" — inilah protokol cas pantas.
  \item \textbf{1. Prinsip Asas Cas Pantas}Kuasa pengecasan = Voltan × Arus (P = U × I). Untuk meningkatkan kuasa, sama ada meningkatkan voltan, meningkatkan arus, atau kedua-duanya. Protokol yang berbeza memilih laluan teknikal yang berbeza.
\end{itemize}

\begin{tcbinfo}{Fungsi "had cas 80\%" iPhone — Tukar peraturan emas ke tetapan sistem}
Siri iPhone 15/16 (iOS 17 ke atas) telah jadikan “peraturuan 20\%-80\%” sebagai fungsi sistem. Selepas aktifkan, telefon cas hingga 80\% akan auto berhenti cas, dengan itu melambatkan penuaan bateri dengan mendadak. Laluan tetapan:

\texttt{Tetapan > Bateri > Kesihatan Bateri \& Cas > Pengoptimuman Cas > Had 80\%}

Tiga pilihan maksud: (1)\textbf{Tidak had}— cas hingga 100\%; (2)\textbf{Optimum cas bateri}— belajar rutin anda, hanya pada bahagian akhir pagi tunda hingga 100\%; (3)\textbf{Had 80\%}— terus berhenti di 80\%, paling jaga bateri. Jangka panjang di rumah/pejabat ada kondisi cas, sangat disyorkan pilih had 80\%; kerap travel, perlu bateri penuh keluar, boleh pilih optimum cas bateri.

\end{tcbinfo}

\begin{tcbinfo}{Pengecas GaN — Rakan kombat travel}
Anda mungkin perhati pengecas pantas sekarang semakin kecil, di sebaliknya adalah\textbf{GaN (Gallium Nitride)}bahan semikonduktor jenis baharu. Berbanding pengecas berasaskan silikon tradisional, pengecas GaN ada beberapa kelebihan:

\begin{itemize}
  \item \textbf{Saiz kecil}: Kuasa sama, pengecas GaN boleh buat 1/2 malah 1/3 saiz silikon tradisional, pengecas 65W GaN lebih kecil daripada pengecas 5W Apple lama.
  \item \textbf{Kecekapan tinggi, haba rendah}: Kehilangan pensuisan gallium nitride jauh lebih rendah daripada silikon, kecekapan penukaran tenaga boleh capai 95\% ke atas, cas tidak mudah panas.
  \item \textbf{Sesuai travel}: Satu pengecas GaN berbilang port 65W atau 100W, boleh serentak cas telefon, tablet, komputer riba, lebih ringan daripada bawa tiga "bata".
  \item \textbf{Cadangan pembelian}: Cari pengecas GaN sokong protokol PD (65W ke atas boleh cas komputer riba), jenama terkenal seperti Baseus, Anker, Nubia, Xiaomi dll boleh. Perhatikan semak kabel data sepadan (kuasa tinggi perlu kabel Type-C dengan cip E-Marker).
\end{itemize}

\end{tcbinfo}

\booksection{Aliran kerja cas pantas:}

Fasa jabat tangan

\booksubsection{Prinsip Asas Cas Pantas}

Fasa rundingan

: Kedua-dua pihak menentukan tahap kuasa tertinggi yang disokong bersama.

\begin{enumerate}
  \item \textbf{Fasa pengecasan}: Pengecas mengeluarkan voltan dan arus mengikut rundingan, telefon pintar melakukan pengurusan kuasa di pihaknya.
  \item \textbf{Fasa pelarasan}: Apabila paras bateri meningkat, telefon pintar akan mengurangkan kuasa pengecasan secara beransur-ansur untuk melindungi bateri.
  \item \textbf{2. Terperinci Protokol Cas Pantas Utama}USB PD (Power Delivery)
  \item \textbf{— Raja Protokol Universal}Ciri
\end{enumerate}

\booksubsection{Terperinci Protokol Cas Pantas Utama}

\textbf{Tahap voltan}: 5V, 9V, 15V, 20V, PD 3.1 juga menyokong 28V, 36V, 48V.

\begin{itemize}
  \item \textbf{FCP (Fast Charge Protocol)}: Semua model iPhone, Google Pixel, sebahagian telefon pintar andalan Android, komputer riba, dan konsol permainan Switch menyokong protokol PD.
  \item \textbf{Kelebihan}: Keserasian universal yang kuat, satu pengecas boleh mengecas pelbagai peranti; menyokong PPS (Power Supply Boleh Atur Cara), boleh melaraskan voltan dan arus dengan tepat.
  \item \textbf{Cip Surge P1}: Memerlukan antara muka Type-C dan cip khusus, kos lebih tinggi.
  \item \textbf{Qualcomm QC (Quick Charge)}— Perintis Kem Android
  \item \textbf{QC 3.0}: Menyokong voltan boleh laras 3.6V-20V dengan langkah 0.2V, maksimum 18W.
\end{itemize}

\textbf{QC 4.0/4+}: Serasi dengan protokol USB PD, menyokong sehingga 27W, menambah teknologi pengecasan dwilaluan.

\begin{itemize}
  \item \textbf{QC 5.0}: Menyokong sehingga 100W+, 0-50\% hanya memerlukan 5 minit, peningkatan besar dalam pengurangan haba.
  \item \textbf{Keserasian}: Kebanyakan telefon pintar dengan pemproses Snapdragon Qualcomm menyokongnya, tetapi setiap jenama biasanya mengutamakan protokol peribadi mereka sendiri.
  \item \textbf{VOOC/Dart}— Wakil Voltan Rendah Arus Tinggi
  \item \textbf{Cip Surge P1}: Voltan rendah arus tinggi (bermula 5V/4A), haba lebih rendah, kecekapan lebih tinggi.
\end{itemize}

\textbf{VOOC 1.0}: 20W (5V/4A), dikeluarkan pada 2014, memulakan era cas pantas.

\begin{itemize}
  \item \textbf{SuperVOOC}: 50W-65W, menggunakan reka bentuk dwi-sel bersiri.
  \item \textbf{SuperVOOC 2.0}: 65W, mengecas penuh bateri 4000mAh dalam 30 minit.
  \item \textbf{150W/240W Super Flash Charge}: Teknologi cas pantas pengeluaran besar terpantas pada masa kini.
  \item \textbf{Ciri}: Memerlukan pengecas dan kabel data khusus (mempunyai pin sentuhan tambahan), tidak serasi dengan protokol lain.
  \item \textbf{Huawei SuperCharge}— Evolusi Cas Pantas Voltan Tinggi
  \item \textbf{FCP (Fast Charge Protocol)}: 9V/2A, 18W, serasi dengan QC 2.0.
\end{itemize}

\textbf{SCP (Super Charge Protocol)}: Menyokong kuasa lebih tinggi, dari 22.5W ke 66W dan kemudian 100W.

\begin{itemize}
  \item \textbf{Ciri teknikal}: Menggunakan teknologi pam cas, menukar voltan tinggi kepada voltan rendah arus tinggi untuk input bateri, kecekapan sehingga 98\%.
  \item \textbf{Keserasian}: Khusus untuk telefon pintar Huawei/Honor, sebahagian pengecas menyokong protokol PD.
  \item \textbf{Protokol}— Kebangkitan Protokol Buatan Sendiri
  \item \textbf{Cip Surge P1}: Cip pengurusan cas pantas buatan Xiaomi, menyokong cas pantas berwayar 120W.
\end{itemize}

\textbf{Cip Surge G1}: Cip pengurusan bateri, bekerjasama dengan P1 untuk mencapai cas pantas yang lebih cekap.

\begin{itemize}
  \item \textbf{Ciri teknikal}: Menggunakan skim dwi-sel bersiri + pam cas, menyokong keserasian pelbagai protokol.
  \item \textbf{Cas Pantas Apple}— Konservatif tetapi Stabil
  \item \textbf{Protokol}: Hanya menyokong protokol USB PD.
\end{itemize}

\textbf{4. Protokol Pengecasan Wayarles}: iPhone 8 dan selepas menyokong cas pantas, maksimum kira-kira 27W (iPhone 15 Pro Max boleh mencapai 27W+).

\begin{itemize}
  \item \textbf{Ciri}: Kuasa agak konservatif, tetapi lengkung pengecasan dioptimumkan dengan baik, kawalan haba sangat cemerlang.
  \item \textbf{Keperluan aksesori}: Memerlukan kabel USB-C ke Lightning (atau kabel USB-C ke USB-C) dan pengecas PD.
  \item \textbf{FCP (Fast Charge Protocol)}Protokol
  \item \textbf{Kuasa Maksimum}Laluan Teknikal
\end{itemize}

\booksubsection{Jadual Perbandingan Protokol Cas Pantas}

\begin{center}
\begin{tabularx}{\linewidth}{|X|X|X|X|}
\hline
\textbf{Ciri} & \textbf{240W} & \textbf{SuperVOOC} & \textbf{Cip Surge P1} \tabularnewline
\hline
QC 5.0 & 100W & Voltan tinggi boleh laras & Peranti Qualcomm \tabularnewline
VOOC/Dart & 240W & Voltan rendah arus tinggi & Khusus keluarga OPPO \tabularnewline
SuperCharge & 100W & Pam cas & Huawei/Honor \tabularnewline
Cas Pantas Apple & \textasciitilde{}27W & Protokol PD & iPhone/iPad \tabularnewline
4. Protokol Pengecasan Wayarles & Pengecasan wayarles juga mempunyai standard protokolnya sendiri: & Standard Qi & : Dirumuskan oleh WPC (Wireless Power Consortium), standard yang paling universal. Qi asas menyokong 5W-15W. \tabularnewline
\hline
\end{tabularx}
\end{center}

\booksubsection{Protokol Pengecasan Wayarles}

: Dikeluarkan pada 2023, memperkenalkan penempatan magnet (serupa dengan Apple MagSafe), menyokong 15W+.

\begin{itemize}
  \item \textbf{Protokol peribadi}: Setiap jenanda mempunyai pengecasan wayarles berkelajuan tinggi sendiri, seperti cas wayarles 50W Xiaomi, cas wayarles 50W Huawei, dll., memerlukan pengecas khusus.
  \item \textbf{5. Panduan Pembelian Cas Pantas}Perhatikan padanan protokol
  \item \textbf{: Apabila membeli pengecas, pastikan ia menyokong protokol cas pantas telefon pintar anda. Pengecas pihak ketiga biasanya akan menyenaraikan protokol yang disokong.}Perhatikan kabel data
\end{itemize}

\booksubsection{Panduan Pembelian Cas Pantas}

\begin{itemize}
  \item \textbf{Perhatikan pensijilan}: Pengguna Apple disyorkan memilih kabel data dengan pensijilan MFi; pengguna Android memilih ori jenama atau pihak ketiga terkenal.
  \item \textbf{Perhatikan haba}: Telefon pintar menjadi panas semasa cas pantas adalah biasa, tetapi jika terlalu panas untuk dipegang, itu tidak normal, disyorkan untuk berhenti mengecas atau menukar pengecas.
  \item \textbf{Petua Keselamatan Cas Pantas}Elakkan mengecas sambil bermain permainan berat, ia akan memburukkan lagi pemanasan.
  \item \textbf{Jangan gunakan pengecas dan kabel data berkualiti rendah, boleh menyebabkan risiko keselamatan.}Buka sarung telefon pintar yang tebal semasa mengecas untuk membantu penyejukan.
\end{itemize}

\begin{tcbwarn}{Jika pengecas atau telefon pintar menjadi panas secara luar biasa atau berbau, hentikan penggunaan dengan serta-merta.}
\begin{itemize}
  \item 4. Ini Barulah Tip Jimat Bateri Sebenar
  \item Skrin (raja penggunaan bateri)
  \item : Kurangkan kecerahan, perpendek masa auto-kunci. Telefon pintar OLED boleh menjimatkan banyak bateri dengan mengaktifkan mod gelap.
  \item Isyarat lemah
\end{itemize}

\end{tcbwarn}

\booksection{IV. Teknologi bateri baharu: Bateri anod silikon-karbon}

Jika anda perhati pengumuman telefon flagship 2024-2026, “anod silikon-karbon” hampir menjadi sinonim bateri premium. Ini adalah penaiktaraf bahan anod bateri litium yang penting, prinsip sebenarnya tidak kompleks:

\begin{itemize}
  \item \textbf{Anod grafit tradisional}: Bateri telefon sentiasa guna grafit sebagai anod, teknologi matang, stabil, tetapi ketumpatan tenaga hampir capai had teori (kira-kira 372mAh/g).
  \item \textbf{Campur bahan silikon}: Kapasiti teori silikon kira-kira 10 kali grafit (kira-kira 4200mAh/g). Campur sedikit silikon ke grafit buat "anod komposit silikon-karbon", boleh tingkatkan ketumpatan tenaga seluruh bateri.
  \item \textbf{Peningkatan ketumpatan tenaga}Isi padu sama, kapasiti bateri anod silikon-karbon boleh tingkat kira-kira\textbf{15\%-20\%}. Ini kunci telefon flagship 2024-2026 boleh muat 5500-6000mAh bateri besar dalam badan 7-9mm.
\end{itemize}

\begin{tcbinfo}{Model mana guna?}
\begin{itemize}
  \item \textbf{Honor}: Siri Magic perintis komersial anod silikon-karbon, Magic6 Pro, Magic7 siri, Magic V3 skrin lipat semuanya guna.
  \item \textbf{OnePlus}: Xiaomi 14 Ultra, siri Xiaomi 15 dilengkapi bateri anod silikon-karbon generasi baharu.
  \item \textbf{vivo}: X100 Ultra, X200 Pro siri guna skim anod silikon-karbon ketumpatan tinggi.
  \item \textbf{Huawei/OPPO dll}: 2025-2026 flagship secara umum ikut, anod silikon-karbon sedang dari "teknologi hitam" jadi standard model premium.
\end{itemize}

\end{tcbinfo}

\textbf{Kesan praktikal kepada pengguna}: Benefit terbesar anod silikon-karbon adalah “kapasiti bateri isi padu sama lebih besar” — telefon boleh buat lebih nipis, pada masa yang sama jangka hayat tidak berkurang; atau badan tidak tebal, tetapi bateri dari 4500mAh naik ke 5500mAh. Maksudnya, anda tidak perlu bayar “tebal berat” untuk jangka hayat. Perhatikan, bateri anod silikon-karbon dalam suhu rendah dan pengurusan cas penuh masih perlu bersama cip cas pantas untuk optimumkan, ini juga salah satu sebab pengeluar cadangkan kekal 80\% cas.

\booksection{V. Ini tip jimat bateri sebenar}

\begin{itemize}
  \item \textbf{5G}: Jika isyarat tidak baik, tukar secara manual kembali ke 4G, ketahanan bateri akan jauh lebih baik.
  \item \textbf{Perkhidmatan lokasi}: Untuk App bukan navigasi, tetapkan kebenaran lokasi kepada "hanya semasa menggunakan".
  \item \textbf{Hadkan aktiviti latar belakang}: iPhone di Tetapan > Umum > Segar Semula App Latar Belakang; Android di tetapan bateri aplikasi.
  \item \textbf{Mari periksa: Lakukan pemeriksaan kesihatan pada bateri anda}Periksa tahap kesihatan: Cari dan catat kapasiti maksimum bateri anda.
  \item \textbf{Analisis butiran penggunaan bateri: Lihat carta dan senarai App dalam tetapan bateri, cari punca utama penggunaan bateri.}Rancang langkah balas: Hadkan App yang menggunakan bateri secara luar biasa di latar belakang.
\end{itemize}

\begin{tcbpractice}{Bab 16: Sistem Akaun \& Perkhidmatan Awan}
\begin{enumerate}
  \item Saya ada seorang kawan, telefon pintarnya dicuri tahun lepas. Perkara pertama yang terlintas di fikirannya bukan "berapa harganya telefon itu", tetapi "adakah foto-foto di dalamnya boleh diambil kembali?" Malangnya, dia tidak mengaktifkan sandaran awan. Telefon pintar adalah harta benda luaran, tetapi data dalam akaun — jika hilang, ia tidak akan kembali selama-lamanya. Bab ini kita akan bincangkan cara memberikan insurans sebenar untuk kehidupan digital anda.
  \item 1. Akaun telefon pintar anda, sebenarnya mengawal berapa banyak perkara?
  \item Ramai orang fikir akaun hanya digunakan untuk memuat turun App. Sebenarnya, perkara yang terikat di sebalik akaun telefon pintar jauh lebih banyak daripada yang anda bayangkan:
\end{enumerate}

\end{tcbpractice}
